% Generated draft: requires mathematical and source review.
\documentclass[11pt]{article}
\usepackage[margin=25mm]{geometry}
\usepackage{fontspec}
\setmainfont{Latin Modern Roman}
\usepackage{amsmath,amssymb,mathtools,unicode-math}
\setmathfont{Latin Modern Math}
\usepackage{enumitem,needspace,xurl,hyperref}
\setlength{\parindent}{0pt}
\setlength{\parskip}{5pt}
\setlength{\emergencystretch}{3em}
\newcommand{\N}{\mathbb{N}}
\newcommand{\Z}{\mathbb{Z}}
\newcommand{\Q}{\mathbb{Q}}
\newcommand{\R}{\mathbb{R}}
\newcommand{\C}{\mathbb{C}}
\newcommand{\F}{\mathbb{F}}
\newcommand{\A}{\mathbb{A}}
\newcommand{\PP}{\mathbb{P}}
\newcommand{\E}{\mathbb{E}}
\newcommand{\eps}{\varepsilon}
\newcommand{\dd}{\,\mathrm{d}}
\begin{document}
\section*{UnsolvedMath ID: 20000734}
\textbf{Source problem code:} catalogue entry
\subsection{Definitions and mathematical statement}
Use the standard Nahm-sum convention
\[
f_{A,B,C}(\tau)=\sum_{n\in\Z_{\geq0}^r}
\frac{q^{\frac12 n^TAn+B^Tn+C}}{\prod_{j=1}^r(q;q)_{n_j}},
\quad q=e^{2\pi i\tau},\quad\Im\tau>0,
\]
where \(A\) is rational symmetric positive definite, \(B\in\Q^r\), \(C\in\Q\), and
\((a;q)_m=\prod_{j=0}^{m-1}(1-aq^j)\), with \((a;q)_\infty\) its convergent infinite-product version. The current record refers to \(f_{A,B,C}\) without reproducing its introductory definition; this convention must be confirmed against the full workshop context.

Problem 2.3 asks for explicit modular examples of rank \(r\geq2\), descriptions by natural objects, and associated combinatorial identities, in analogy with the two rank-one Rogers--Ramanujan examples mentioned in the source. It does not define natural, prescribe a modular group/weight/multiplier, or ask for an exhaustive classification. Appropriate modular or partition interpretations are therefore part of the requested development.

The attached aid treats the genuinely coupled rank-two matrix
\[
A=\begin{pmatrix}4&1\\1&1\end{pmatrix},\qquad
B_\varepsilon=(2\varepsilon,1/2),\qquad
C_0=1/120,\quad C_1=49/120,\quad\varepsilon\in\{0,1\}.
\]
Define
\[
G(Q)=\sum_{i\geq0}\frac{Q^{i^2}}{(Q;Q)_i},\qquad
H(Q)=\sum_{i\geq0}\frac{Q^{i(i+1)}}{(Q;Q)_i}.
\]
In the stated convention, the supplied product reductions are expressed with the normalization kept explicit:
\[
q^{-C_0}f_{A,B_0,C_0}=(-q;q)_\infty G(q^2),\qquad
q^{-C_1}f_{A,B_1,C_1}=(-q;q)_\infty H(q^2).
\]
The aid reports the normalized \(f\)'s as eta/Siegel modular-unit products, without specifying the exact group or multiplier. The fractional \(q^{C_\varepsilon}\) prefactors must not be removed when asserting modularity, even though they are removed for ordinary partition generating functions.

The aid's combinatorial interpretations say the first stripped series counts partitions avoiding residues \(0,4,6\pmod{10}\), and the second avoids \(0,2,8\pmod{10}\). It additionally claims a weight- and charge-preserving bijection \(\Phi_\varepsilon\) between:
\begin{itemize}
\item triples \((i,\alpha,\beta)\), with \(i\geq0\), \(\alpha\) a partition with parts at most \(i\), and \(\beta\) a distinct-part partition with parts greater than \(i\), carrying total weight \(2i^2+2\varepsilon i+|\alpha|+|\beta|\);
\item pairs \((\delta,\rho)\), where \(\delta\) is a distinct-part partition and \(\rho\) has even parts separated by gaps at least four, with minimum part at least \(2+2\varepsilon\), carrying weight \(|\delta|+|\rho|\).
\end{itemize}
The charge condition is that \(i\) equals the number of parts of \(\rho\); empty partitions are allowed. The map itself is not supplied in the record, so this is a reported bijection claim rather than an explicit construction here. The aid distinguishes recovering published identities from its claimed additional bijection and does not solve the broader higher-rank problem. All such partial results remain unreviewed supplied claims.
\subsection{Short English statement}
Find natural modular and combinatorial descriptions of explicit rank-at-least-two Nahm sums. The supplied aid gives two coupled rank-two examples, preserving their fractional modular normalizations and distinguishing product identities from an additional claimed partition bijection.
\subsection{Sources}
\begin{itemize}
\item AIM Problem 2.3, workshop \emph{Mock modular forms in combinatorics and arithmetic geometry}, 8--12 March 2010 according to the supplied verification: \url{https://aimath.org/WWN/mockmodular/mockmodular.pdf}.
\item Local registry \texttt{data/unsolvedmath/problems.json}, permanent ID \(20000734\), exact source code \texttt{catalogue entry}; canonical \texttt{aim-analytic-number-theory-notes.json}, notes[97].
\item Related local Problems 2.1 and 2.2, permanent IDs \(20000732\) and \(20000733\), source codes \texttt{AIM-ANALYTIC\_NUMBER\_THEORY-0096} and \texttt{-0097}, canonical notes[95] and notes[96], discuss the same parameters but also omit the full introductory definition.
\item The source names Rogers--Ramanujan identities. The attached aid calls the coupled product identities published but supplies no article citations or explicit \(\Phi_\varepsilon\). Attempt 1 is explicitly machine-generated, not peer reviewed, and requires independent verification before citation.
\end{itemize}
\end{document}
